% !tp-example Basic usage
\documentclass{article}
\begin{document}
\wcTpNotebook{
  \wcScript[type=tp/markdown]|
    # Notebook example
    Here is a notebook that combines AsciiDoc, HTML, LaTeX, Markdown and reStructuredText as markup languages, and Python and JavaScript as programming languages.

    ## Calculation
    We perform a simple multiplication in Python to demonstrate the calculation.
  |
  \wcScript[type=tp/python]|
    x = 6 * 7
    print(x)
  |
  \wcScript[type=tp/restructuredtext,doctest]|
    .. h2::

       Verification

    We use the doctest blocks in reStructuredText to verify the output of the Python code.

    >>> print(6 * 7)
    42
  |
  \wcScript[type=tp/html]|
    In JavaScript, we would have written the following code:
  |
  \wcScript[type=tp/javascript]|
    const x = 6 * 7;
    console.log(x)
  |
  \wcScript[type=tp/asciidoc]|
    == Result
    The result of the multiplication is displayed using Asciidoc and Mathjax:

    latexmath:[x = 6 * 7 = 42]
  |
  \wcScript[type=tp/latex]|
    \section*{LaTeX and TP styles}
    The same result can be written in LaTeX: $6 \times 7 = 42$.

    \begin{itemize}
      \item Ordinary lists use the shared TP styles.
      \item Nested lists keep their hierarchy.
        \begin{itemize}
          \item This is a nested item.
        \end{itemize}
    \end{itemize}

    \wcTpCallout[variant=info,heading=LaTeX inside a component]{
      \textbf{The component controls its appearance.}
      \begin{itemize}
        \item Keep the LaTeX source readable.
        \item Reuse the same TP typography and list styles.
      \end{itemize}
    }
  |
}
\end{document}
